% !TeX root = part3.tex
\documentclass[../final.tex]{subfiles}

\begin{document}

\ifSubfilesClassLoaded{\setcounter{chapter}{2}\setcounter{section}{0}}{}

% Источники: фотографии 1.jpg--13.jpg, предоставленные пользователем.
% На рукописной странице указана дата 13.09, без года.
% Повторяющиеся записи тетради и слайдов объединены по темам.
\chapter{Многоэлектронные атомы}
\rsubtitle{Атом гелия. Электронные конфигурации и термы}

\section{Связь с одноэлектронной моделью}

В предыдущей лекции энергия внешнего электрона щелочного атома
описывалась через \rconcept{квантовый дефект} \(\Delta_\ell\):
\begin{equation*}
    E_{n\ell}=-\frac{Z_a^2Ry}{(n-\Delta_\ell)^2}.
\end{equation*}
Здесь \(Z_a\) --- заряд атомного остова, а \(Ry\) --- энергия Ридберга.
Для нейтрального щелочного атома \(Z_a=1\).

Поправка тонкой структуры водородоподобного уровня имеет вид
\begin{equation*}
    \delta E_{nj}
    =-\frac{\alpha^2 Z^4Ry}{n^3}
    \left(\frac{1}{j+1/2}-\frac{3}{4n}\right),
\end{equation*}
где \(\alpha\) --- постоянная тонкой структуры.

Для внутренних оболочек используется модель экранированного заряда:
\begin{equation*}
    E_n=-\frac{Ry(Z-a)^2}{n^2}.
\end{equation*}
При переходе с верхнего уровня \(n_i\) на нижний \(n_f\)
энергия испущенного фотона равна
\begin{equation*}
    \hbar\omega=E_i-E_f
    =Ry(Z-a)^2\left(\frac{1}{n_f^2}-\frac{1}{n_i^2}\right).
\end{equation*}
Такая запись предполагает одинаковую эффективную поправку \(a\)
для двух уровней; при разных поправках берётся разность соответствующих
энергий. Далее рассмотрим явный учёт взаимодействия электронов.

\section{Атом гелия в теории возмущений}

\subsection{Основное состояние}

В атоме гелия два электрона и ядро с зарядом \(Z=2\).
В \rassumption{системе СГС, при неподвижном ядре} гамильтониан равен
\begin{align*}
    \hat H&=\hat H_0+\hat V,\\
    \hat H_0&=\sum_{i=1}^{2}
    \left(\frac{\hat{\mathbf p}_i^{\,2}}{2m}-\frac{Ze^2}{r_i}\right),
    &\hat V&=\frac{e^2}{r_{12}},
    \qquad r_{12}=|\mathbf r_1-\mathbf r_2|.
\end{align*}
В нулевом приближении электроны не взаимодействуют. Для конфигурации
\(1s^2\) пространственная волновая функция и энергия имеют вид
\begin{equation*}
    \Phi_0(\mathbf r_1,\mathbf r_2)
    =\psi_{1s}(\mathbf r_1)\psi_{1s}(\mathbf r_2),
    \qquad E_0^{(0)}=2E_{1s}.
\end{equation*}
Поправка первого порядка:
\begin{equation*}
    \Delta E_0^{(1)}
    =\iint \psi_{1s}^{*}(\mathbf r_1)\psi_{1s}^{*}(\mathbf r_2)
    \frac{e^2}{r_{12}}
    \psi_{1s}(\mathbf r_1)\psi_{1s}(\mathbf r_2)
    \,d^3r_1\,d^3r_2.
\end{equation*}
Звёздочка означает комплексное сопряжение. Спиновая часть здесь
не выписана; для основного состояния она антисимметрична.

\subsection{Возбуждённое состояние и обменное вырождение}

Пусть один электрон находится в состоянии \(a=1s\), другой ---
в состоянии \(b=2s\). Тогда
\begin{equation*}
    E_{\text{возб}}^{(0)}=E_{1s}+E_{2s}.
\end{equation*}
Одна и та же энергия соответствует двум пространственным функциям:
\begin{align*}
    \Phi_1&=\psi_a(\mathbf r_1)\psi_b(\mathbf r_2),\\
    \Phi_2&=\psi_b(\mathbf r_1)\psi_a(\mathbf r_2).
\end{align*}
Поэтому используем теорию возмущений для вырожденного уровня.
Введём \rconcept{кулоновский интеграл}
\begin{equation*}
    C=\iint |\psi_a(\mathbf r_1)|^2
    \frac{e^2}{r_{12}}|\psi_b(\mathbf r_2)|^2
    \,d^3r_1\,d^3r_2
\end{equation*}
и \rconcept{обменный интеграл}
\begin{equation*}
    A=\iint \psi_a^{*}(\mathbf r_1)\psi_b^{*}(\mathbf r_2)
    \frac{e^2}{r_{12}}
    \psi_b(\mathbf r_1)\psi_a(\mathbf r_2)
    \,d^3r_1\,d^3r_2.
\end{equation*}
В рассматриваемом случае \(A\) вещественен. Матрица возмущения
в базисе \(\Phi_1,\Phi_2\) равна
\begin{equation*}
    V=\begin{pmatrix}C&A\\A&C\end{pmatrix}.
\end{equation*}
Секулярное уравнение даёт
\begin{equation*}
    \begin{vmatrix}C-\lambda&A\\A&C-\lambda\end{vmatrix}
    =(C-\lambda)^2-A^2=0.
\end{equation*}
Следовательно,
\begin{equation*}
    \rboxed{\Delta E_{\pm}^{(1)}=C\pm A},
    \qquad
    \Phi_\pm=\frac{\Phi_1\pm\Phi_2}{\sqrt2}.
\end{equation*}
Функция \(\Phi_+\) симметрична по перестановке электронов,
а \(\Phi_-\) антисимметрична. Нормировка \(1/\sqrt2\) относится
к различным ортонормированным орбиталям \(a\ne b\).

\subsection{Поправка к энергии}

Для произвольной, в том числе ненормированной, функции
формула первого порядка записывается как
\begin{equation*}
    \Delta W=
    \frac{\int\Psi^{*}\hat V\Psi\,d\tau}
         {\int\Psi^{*}\Psi\,d\tau}.
\end{equation*}
После суммирования по спиновым переменным для \(a\ne b\) получаем
\begin{align*}
    \Delta W_\pm
    &=\frac12\int
    \bigl[\psi_a(\mathbf r_1)\psi_b(\mathbf r_2)
    \pm\psi_b(\mathbf r_1)\psi_a(\mathbf r_2)\bigr]^{*}\\
    &\qquad\times\frac{e^2}{r_{12}}
    \bigl[\psi_a(\mathbf r_1)\psi_b(\mathbf r_2)
    \pm\psi_b(\mathbf r_1)\psi_a(\mathbf r_2)\bigr],d\tau\\
    &=C\pm A,
    \qquad d\tau=d^3r_1\,d^3r_2.
\end{align*}
Эквивалентная симметризованная запись обменного интеграла со слайда:
\begin{align*}
    A=\frac12\int\Bigl[
    &\psi_a^{*}(\mathbf r_1)\psi_b^{*}(\mathbf r_2)
     \psi_b(\mathbf r_1)\psi_a(\mathbf r_2)\\
    +{}&\psi_b^{*}(\mathbf r_1)\psi_a^{*}(\mathbf r_2)
     \psi_a(\mathbf r_1)\psi_b(\mathbf r_2)
    \Bigr]\frac{e^2}{r_{12}}\,d\tau.
\end{align*}
Для одинаковых орбиталей \(a=b\) пространственная функция равна
\(\psi_a(\mathbf r_1)\psi_a(\mathbf r_2)\), и поправка составляет \(C\).
Таким образом, полная энергия в первом порядке:
\begin{equation*}
    \rboxed{
    W=\begin{cases}
        W_a+W_b+C\pm A,&a\ne b,\\
        2W_a+C,&a=b.
    \end{cases}}
\end{equation*}

\section{Принцип Паули и спиновые функции}

\subsection{Антисимметрия полной волновой функции}

Электрон имеет спин \(s=1/2\). Согласно \rconcept{принципу Паули},
полная волновая функция электронов антисимметрична относительно
перестановки любых двух частиц:
\begin{equation*}
    \Psi(\ldots,x_i,\ldots,x_k,\ldots)
    =-\Psi(\ldots,x_k,\ldots,x_i,\ldots).
\end{equation*}
Здесь \(x_i=(\mathbf r_i,\sigma_i)\) включает пространственные
и спиновые переменные. Антисимметрия относится именно к полной функции.

\begin{rbox}{[}
    \rresult{Никакие два электрона не могут находиться в одном
    и том же одноэлектронном квантовом состоянии.}
\end{rbox}

Базисные спиноры обозначим
\begin{equation*}
    \alpha=\begin{pmatrix}1\\0\end{pmatrix},
    \quad m_s=+\frac12;
    \qquad
    \beta=\begin{pmatrix}0\\1\end{pmatrix},
    \quad m_s=-\frac12.
\end{equation*}
В обозначениях со стрелками \(\alpha=\uparrow\), \(\beta=\downarrow\).

\subsection{Парагелий и ортогелий}

Для двух электронов возможны полный спин \(S=0\) и \(S=1\).
Обозначим его проекцию через \(M_S\), чтобы отличать её от проекции
\(m_s\) отдельного электрона.

\rconcept{Парагелий} соответствует синглетному состоянию
\(S=0\), \(M_S=0\). Его спиновая функция антисимметрична:
\begin{equation*}
    \chi_{00}=\frac{\alpha(1)\beta(2)-\beta(1)\alpha(2)}{\sqrt2}.
\end{equation*}
При \(a\ne b\) полная функция имеет вид
\begin{align*}
    \Psi_{00}
    &=\Phi_+\chi_{00}\\
    &=\frac12\bigl[\psi_a(\mathbf r_1)\psi_b(\mathbf r_2)
       +\psi_b(\mathbf r_1)\psi_a(\mathbf r_2)\bigr]
       \bigl[\uparrow\downarrow-\downarrow\uparrow\bigr].
\end{align*}
Ей соответствует энергия \(W_a+W_b+C+A\).

\rconcept{Ортогелий} соответствует триплетному состоянию \(S=1\).
Все три спиновые функции симметричны:
\begin{align*}
    \chi_{11}&=\alpha(1)\alpha(2),\\
    \chi_{10}&=\frac{\alpha(1)\beta(2)+\beta(1)\alpha(2)}{\sqrt2},\\
    \chi_{1,-1}&=\beta(1)\beta(2).
\end{align*}
Поэтому пространственная часть должна быть антисимметричной:
\begin{align*}
    \Psi_{11}
    &=\frac{1}{\sqrt2}
      \bigl[\psi_a(\mathbf r_1)\psi_b(\mathbf r_2)
      -\psi_b(\mathbf r_1)\psi_a(\mathbf r_2)\bigr]
      \uparrow\uparrow,\\
    \Psi_{10}
    &=\frac12
      \bigl[\psi_a(\mathbf r_1)\psi_b(\mathbf r_2)
      -\psi_b(\mathbf r_1)\psi_a(\mathbf r_2)\bigr]
      \bigl[\uparrow\downarrow+\downarrow\uparrow\bigr],\\
    \Psi_{1,-1}
    &=\frac{1}{\sqrt2}
      \bigl[\psi_a(\mathbf r_1)\psi_b(\mathbf r_2)
      -\psi_b(\mathbf r_1)\psi_a(\mathbf r_2)\bigr]
      \downarrow\downarrow.
\end{align*}
Энергия триплета равна \(W_a+W_b+C-A\).
Обменное расщепление между синглетом и триплетом составляет
\begin{equation*}
    \rboxed{W_{\text{пара}}-W_{\text{орто}}=2A}.
\end{equation*}
При \(a=b\) антисимметричная пространственная комбинация обращается
в нуль. Поэтому конфигурация \(1s^2\) допускает только синглет:
\begin{equation*}
    \Psi_{\text{осн}}
    =\psi_{1s}(\mathbf r_1)\psi_{1s}(\mathbf r_2)\chi_{00}.
\end{equation*}

\subsection{Диаграмма Гротриана}

Уровни гелия образуют синглетную и триплетную системы.
В обозначении терма \({}^{2S+1}L\) верхний индекс задаёт
мультиплетность, а буквы \(S,P,D,F\) соответствуют
\(L=0,1,2,3\).

\begin{rbox}[left=12pt,right=12pt,top=12pt,bottom=10pt]{[]}
    \centering
    \begin{tikzpicture}[x=0.95cm,y=0.25cm,>=Stealth,
        every node/.style={font=\small},line width=0.7pt]
        \draw[->] (0,0)--(0,27) node[above] {\(E_{\text{возб}},\ \text{эВ}\)};
        \foreach \v in {0,5,10,15,20} {
            \draw (-0.10,\v)--(0.10,\v);
            \node[left] at (-0.1,\v) {\(\v\)};
        }
        \draw[densely dotted] (0,24.6)--(13.1,24.6);
        \node[left] at (-0.1,24.6) {\(24.6\)};
        \node at (3.4,29) {Парагелий: \(S=0\)};
        \node at (10.25,29) {Ортогелий: \(S=1\)};
        \foreach \x/\term in {1.3/S,2.9/P,4.5/D,6.1/F} {
            \node at (\x,26.5) {\({}^{1}\term\)};
        }
        \foreach \x/\term in {8.1/S,9.7/P,11.3/D,12.9/F} {
            \node at (\x,26.5) {\({}^{3}\term\)};
        }
        % Уровни размещены схематически: важны серии и переходы.
        \draw (0.95,0)--(1.65,0) node[right] {\(1s^2\)};
        \foreach \x/\y in {1.3/20.6,2.9/21.3,8.1/19.8,9.7/20.5,
            1.3/22.65,2.9/23.15,4.5/23.1,8.1/22.4,9.7/22.95,11.3/23.05,
            1.3/23.6,2.9/23.85,4.5/23.8,6.1/23.8,
            8.1/23.5,9.7/23.8,11.3/23.8,12.9/23.8} {
            \draw (\x-0.35,\y)--(\x+0.35,\y);
        }
        \foreach \x in {1.3,2.9,4.5,6.1,8.1,9.7,11.3,12.9} {
            \foreach \y in {24.15,24.35,24.5} {
                \draw (\x-0.35,\y)--(\x+0.35,\y);
            }
        }
        \draw[->] (2.9,21.3)--(1.3,0.2);
        \draw[->] (2.9,23.15)--(1.3,20.6);
        \draw[->] (1.3,23.6)--(2.9,21.3);
        \draw[->] (4.5,23.1)--(2.9,21.3);
        \draw[->] (6.1,23.8)--(4.5,23.1);
        \draw[->] (2.9,23.85)--(4.5,23.1);
        \draw[->] (9.7,22.95)--(8.1,19.8);
        \draw[->] (8.1,23.5)--(9.7,20.5);
        \draw[->] (11.3,23.05)--(9.7,20.5);
        \draw[->] (12.9,23.8)--(11.3,23.05);
        \draw[->] (9.7,23.8)--(11.3,23.05);
        \draw[rresultcolor,dashed] (1.65,20.6)--(7,20.6);
        \draw[rresultcolor,dashed] (7,19.8)--(8.45,19.8);
        \draw[rresultcolor,<->] (7,19.8)--(7,20.6);
        \node[rresultcolor,below,align=center] at (6.2,18.8)
            {\(2A\simeq0.75\,\text{эВ}\)\\для \(1s2s\)};
        \node[align=left] at (8,8)
            {Энергия зависит от \(S\) и \(L\).\\
             До учёта тонкой структуры\\
             сохраняется вырождение по \(J\).};
    \end{tikzpicture}
    \captionof{figure}{Диаграмма Гротриана для гелия по слайду.
    Положение возбуждённых уровней и стрелки показаны схематически;
    шкала энергии отсчитывается от основного состояния.}
    \label{fig:helium_grotrian}
\end{rbox}

На слайде для \(a=1s\), \(b=2s\) приведены значения
\begin{align*}
    W_a+W_b&=\left(-Z^2-\frac{Z^2}{4}\right)Ry
             =-67.99\,\text{эВ},\\
    C&=10.07\,\text{эВ},\qquad A=1.9\,\text{эВ}.
\end{align*}
\remark[Замечание к исходному слайду.]
Число \(A=1.9\,\text{эВ}\) и подпись \(2A=0.75\,\text{эВ}\)
на диаграмме не согласуются как значения одного и того же интеграла:
первое дало бы \(2A=3.8\,\text{эВ}\).
Оба значения сохранены из источника; связывать их одним численным
расчётом без уточнения используемого приближения нельзя.

\section{Приближение независимых электронов}

В многоэлектронном атоме заменим действие остальных электронов
некоторым средним полем. Тогда гамильтониан представляется суммой
одноэлектронных операторов:
\begin{equation*}
    \hat H_0=\sum_{i=1}^{N}\hat h_i,
    \qquad
    \hat h_i=\frac{\hat{\mathbf p}_i^{\,2}}{2m}
             -\frac{Ze^2}{r_i}+U_{\text{ср}}(\mathbf r_i).
\end{equation*}
Волновую функцию сначала ищем как произведение
одноэлектронных функций --- \rconcept{произведение Хартри}:
\begin{equation*}
    \Phi(\mathbf r_1,\ldots,\mathbf r_N)
    =\varphi_1(\mathbf r_1)\varphi_2(\mathbf r_2)
     \cdots\varphi_N(\mathbf r_N).
\end{equation*}
Уравнение \(\hat H_0\Phi=W_0\Phi\) распадается на систему
\begin{equation*}
    \hat h_i\varphi_i(\mathbf r_i)=W_i\varphi_i(\mathbf r_i),
    \qquad i=1,2,\ldots,N,
\end{equation*}
причём
\begin{equation*}
    \rboxed{W_0=\sum_{i=1}^{N}W_i}.
\end{equation*}
Это энергия заданного нулевого гамильтониана \(\hat H_0\).
Само произведение Хартри ещё не обеспечивает антисимметрию
полной электронной функции.

В центральном поле пространственная часть одноэлектронной функции
разделяется на радиальную и угловую:
\begin{equation*}
    \varphi_i(\mathbf r_i)
    =R_{n_i\ell_i}(r_i)Y_{\ell_i m_{\ell i}}(\theta_i,\phi_i).
\end{equation*}

\subsection{Детерминант Слейтера}

Антисимметричную функцию строят из одноэлектронных
\rconcept{спин-орбиталей} \(\varphi_{p}(x)\), включающих спин.
Для двух электронов
\begin{align*}
    \Psi(x_1,x_2)
    &=\frac{1}{\sqrt2}
    \bigl[\varphi_{p_1}(x_1)\varphi_{p_2}(x_2)
          -\varphi_{p_2}(x_1)\varphi_{p_1}(x_2)\bigr]\\
    &=\frac{1}{\sqrt2}
    \begin{vmatrix}
        \varphi_{p_1}(x_1)&\varphi_{p_2}(x_1)\\
        \varphi_{p_1}(x_2)&\varphi_{p_2}(x_2)
    \end{vmatrix}.
\end{align*}
Для \(N\) электронов получаем \rconcept{детерминант Слейтера}:
\begin{equation*}
    \rboxed{
    \Psi(x_1,\ldots,x_N)=\frac{1}{\sqrt{N!}}
    \begin{vmatrix}
        \varphi_{p_1}(x_1)&\cdots&\varphi_{p_N}(x_1)\\
        \vdots&\ddots&\vdots\\
        \varphi_{p_1}(x_N)&\cdots&\varphi_{p_N}(x_N)
    \end{vmatrix}}.
\end{equation*}
Множитель \(1/\sqrt{N!}\) нормирует функцию при ортонормированных
спин-орбиталях. Перестановка электронов меняет местами строки
и меняет знак определителя. Если две спин-орбитали совпадают,
два столбца одинаковы, определитель равен нулю: так реализуется
принцип исключения Паули.

\section{Электронные конфигурации и оболочки}

\subsection{Квантовые числа электрона}

В одноэлектронном приближении состояние электрона задаётся
четырьмя квантовыми числами. Используются два представления:
\begin{equation*}
    \psi_{n,\ell,m_\ell,m_s}
    \qquad\text{и}\qquad
    \psi_{n,\ell,j,m_j}.
\end{equation*}
Первое соответствует раздельным проекциям орбитального момента
и спина, второе --- полному моменту \(\mathbf j=\boldsymbol\ell+\mathbf s\).

\begin{center}
\renewcommand{\arraystretch}{1.5}
\begin{tabularx}{\linewidth}{|X|l|}
    \hline
    \rconcept{\textbf{Квантовое число}} &
    \rconcept{\textbf{Допустимые значения}}\\
    \hline
    \(n\) --- главное & \(n=1,2,3,\ldots\)\\
    \hline
    \(\ell\) --- орбитальное & \(\ell=0,1,\ldots,n-1\)\\
    \hline
    \(m_\ell\) --- магнитное орбитальное;
    проекция орбитального момента на ось \(z\) в единицах \(\hbar\)
    & \(m_\ell=-\ell,-\ell+1,\ldots,\ell\)\\
    \hline
    \(m_s\) --- магнитное спиновое;
    проекция спина на ось \(z\) в единицах \(\hbar\), \(s=1/2\)
    & \(m_s=\pm1/2\)\\
    \hline
\end{tabularx}
\end{center}

\subsection{Оболочки и подоболочки}

\rconcept{Оболочка} объединяет состояния с одинаковым \(n\),
а \rconcept{подоболочка} --- с одинаковыми \(n\) и \(\ell\).

\begin{center}
\renewcommand{\arraystretch}{1.35}
\begin{tabular}{l|ccccc}
    \hline
    \(n\)&1&2&3&4&5\\
    Обозначение оболочки&\(K\)&\(L\)&\(M\)&\(N\)&\(O\)\\
    \hline
\end{tabular}

\medskip
\begin{tabular}{l|cccc}
    \hline
    \(\ell\)&0&1&2&3\\
    Обозначение подоболочки&\(s\)&\(p\)&\(d\)&\(f\)\\
    \hline
\end{tabular}
\end{center}

Для заданного \(n\)
\begin{equation*}
    \ell=n-1,n-2,\ldots,0,
    \qquad m_\ell=-\ell,\ldots,0,\ldots,\ell.
\end{equation*}
В \rconcept{электронной конфигурации} верхний индекс у обозначения
подоболочки указывает число электронов. Например,
\begin{equation*}
    \mathrm{Na}:\quad 1s^2\,2s^2\,2p^6\,3s^1.
\end{equation*}

\subsection{Заполнение электронных состояний}

Описание электронных оболочек и периодической системы опирается
на три положения:
\begin{enumerate}
    \item Одноэлектронное приближение: движение в эффективном поле
    позволяет выделить оболочки и подоболочки.
    \item Принцип Паули: в состоянии с заданным набором
    \((n,\ell,m_\ell,m_s)\) может находиться только один электрон.
    \item Принцип минимума энергии: основное состояние соответствует
    наименьшей энергии при соблюдении принципа Паули.
\end{enumerate}
Эффективное центральное поле характеризуется зарядом
\(Z_{\text{эф}}(r)\). Потенциальная энергия притяжения электрона равна
\begin{equation*}
    U_{\text{эф}}(r)=-\frac{Z_{\text{эф}}(r)e^2}{r}.
\end{equation*}
\rnote{На слайде записана положительная величина
\(Z_{\text{эф}}(r)e^2/r\); для потенциальной энергии электрона
здесь явно указан знак притяжения.}

Приведённые на слайде качественные тенденции уменьшения энергии:
\begin{equation*}
    Z\uparrow,\qquad n\downarrow,\qquad\ell\downarrow.
\end{equation*}
Сравнение по \(\ell\) относится к уровням при фиксированном \(n\)
в экранированном поле; для чисто кулоновского поля нерелятивистская
энергия от \(\ell\) не зависит.

\section{LS-связь и электронные термы}

\subsection{Приближение Рассела--Саундерса}

В приближении \rconcept{\(LS\)-связи} сначала складываются
орбитальные моменты электронов и отдельно их спины:
\begin{equation*}
    \mathbf L=\sum_{i=1}^{N}\boldsymbol\ell_i,
    \qquad \mathbf S=\sum_{i=1}^{N}\mathbf s_i,
    \qquad \mathbf J=\mathbf L+\mathbf S.
\end{equation*}
Это \rassumption{векторное сложение моментов}; квантовые числа
\(L\) и \(S\) не равны в общем случае суммам чисел \(\ell_i\) и \(s_i\).

В рамках выбранной электронной конфигурации \(\{n_i,\ell_i\}\)
остаются хорошими квантовыми числами. Индивидуальные проекции
\(m_{\ell i}\) и \(m_{s i}\) уже не определяют стационарное состояние.
Без учёта спин-орбитального взаимодействия хорошими числами являются
\begin{equation*}
    L,\quad S,\quad M_L,\quad M_S.
\end{equation*}
Таким образом, состояния характеризуются набором
\begin{equation*}
    \{n_i,\ell_i\},\quad S,L,M_S,M_L.
\end{equation*}
\rconcept{Электронный терм} \({}^{2S+1}L\) объединяет состояния
с заданными конфигурацией, \(S\) и \(L\), но разными проекциями.
До учёта тонкой структуры число этих состояний равно
\begin{equation*}
    g=(2S+1)(2L+1).
\end{equation*}
При учёте полного момента используются обозначения
\({}^{2S+1}L_J\), где
\begin{equation*}
    J=|L-S|,|L-S|+1,\ldots,L+S.
\end{equation*}

\subsection{Пример: конфигурация \(1s^2\,2s^2\,2p^2\)}

Замкнутые подоболочки \(1s^2\) и \(2s^2\) дают нулевые суммарные
моменты. Термы определяются двумя электронами подоболочки \(2p\):
\begin{equation*}
    \ell_1=\ell_2=1,\qquad
    m_{\ell 1},m_{\ell 2}=-1,0,1,
    \qquad m_{s1},m_{s2}=\pm\frac12.
\end{equation*}
В подоболочке имеется шесть спин-орбиталей, поэтому число допустимых
двухэлектронных состояний равно
\begin{equation*}
    \binom62=15.
\end{equation*}
Подсчёт состояний с заданными
\(M_L=m_{\ell1}+m_{\ell2}\) и \(M_S=m_{s1}+m_{s2}\)
даёт таблицу, соответствующую разбору на рукописной странице:
\begin{center}
\renewcommand{\arraystretch}{1.3}
\begin{tabular}{c|ccc}
    \hline
    \(M_L\backslash M_S\)&\(1\)&\(0\)&\(-1\)\\
    \hline
    \(2\)&0&1&0\\
    \(1\)&1&2&1\\
    \(0\)&1&3&1\\
    \(-1\)&1&2&1\\
    \(-2\)&0&1&0\\
    \hline
\end{tabular}
\end{center}
Например, при \(M_L=2\) оба электрона имеют \(m_\ell=1\).
Их спины обязаны быть противоположны, поэтому \(M_S=0\).
Состояние \(M_L=2,\ M_S=1\) запрещено принципом Паули.

Состояние с максимальным \(M_L=2\) выделяет терм
\(L=2,S=0\), то есть \({}^{1}D\), содержащий пять состояний.
После их выделения остаётся терм \(L=1,S=1\), то есть
\({}^{3}P\), с девятью состояниями. Последнее состояние имеет
\(L=0,S=0\) и образует терм \({}^{1}S\).

\begin{center}
\renewcommand{\arraystretch}{1.3}
\begin{tabular}{c|c|c|c}
    \hline
    \rconcept{Терм}&\(L\)&\(S\)&\((2L+1)(2S+1)\)\\
    \hline
    \({}^{1}D\)&2&0&5\\
    \({}^{3}P\)&1&1&9\\
    \({}^{1}S\)&0&0&1\\
    \hline
\end{tabular}
\end{center}
Итого \(5+9+1=15\): все допустимые состояния учтены.

\section{Тонкая структура}

\subsection{Водородоподобный атом}

С учётом поправки тонкой структуры энергия уровня записывается как
\begin{equation*}
    \rboxed{
    W_{nj}=-\frac{RyZ^2}{n^2}
    \left[1+\frac{\alpha^2Z^2}{n}
    \left(\frac{1}{j+1/2}-\frac{3}{4n}\right)\right]}.
\end{equation*}
На слайде полный момент одного электрона обозначен через \(J\);
здесь для него используется \(j\), а \(J\) обозначает момент
всего многоэлектронного атома.

\subsection{Атом с несколькими электронами}

Спин-орбитальная поправка в приближении \(LS\)-связи имеет вид
\begin{equation*}
    \Delta\hat W'
    =\sum_{i=1}^{N}a_i\,\mathbf s_i\cdot\boldsymbol\ell_i
    \simeq A_{LS}\,\mathbf S\cdot\mathbf L.
\end{equation*}
Здесь \rassumption{операторы моментов выражены в единицах \(\hbar\)}.
Коэффициент \(A_{LS}\) имеет размерность энергии. На слайде он
обозначен буквой \(A\); индекс добавлен, чтобы отличать его от
обменного интеграла гелия.

Из равенства \(\mathbf J=\mathbf L+\mathbf S\) следует
\begin{equation*}
    \mathbf L\cdot\mathbf S
    =\frac12\left(\mathbf J^2-\mathbf L^2-\mathbf S^2\right).
\end{equation*}
Поэтому в состоянии с определёнными \(L,S,J\)
\begin{equation*}
    \left\langle\mathbf L\cdot\mathbf S\right\rangle
    =\frac12\bigl[J(J+1)-L(L+1)-S(S+1)\bigr].
\end{equation*}
Поправка к энергии терма равна
\begin{equation*}
    \rboxed{\Delta W'_J=
    \frac{A_{LS}}{2}
    \bigl[J(J+1)-L(L+1)-S(S+1)\bigr]}.
\end{equation*}
Таким образом, спин-орбитальное взаимодействие снимает вырождение
по \(J\) и расщепляет терм на уровни \({}^{2S+1}L_J\).

\end{document}
